\begin{table*}
\footnotesize
\center
\begin{tabular}{@{}lcccccc@{}}
\toprule
\textbf{Platform}        & \textbf{Multi-Turn Chat} & \textbf{Partial Obs} & \textbf{Long Horizon} & \textbf{Embodied} & \textbf{Quantitative}\\ 
\midrule
Overcooked \cite{carroll2019utility} & \multicolumn{1}{c}{{\color[HTML]{009901} }}  &   & & \checkmark & \checkmark   \\
CerealBar \cite{suhr-etal-2019-executing}    & \checkmark  & \checkmark  & & \checkmark & \checkmark                                                                     \\
Habitat AI \cite{savva2019habitat}   &  & \checkmark   &  & \checkmark & \checkmark \\
LLM-Coord \cite{agashe2023llm} & \checkmark & \checkmark & & & \checkmark\\
Generative Agents \cite{park2023generative} & \checkmark & \checkmark & \checkmark &  & \\
PARTNR \cite{chang2024partnr}   &  & \checkmark   &  & \checkmark & \checkmark \\
MineLand \cite{yu2024mineland}   & \checkmark & \checkmark  &  & \checkmark &  \\
\midrule
\textbf{\mindcraft{} (ours)}     & \textbf{\checkmark} & \checkmark & \checkmark & \checkmark & \checkmark \\ \bottomrule
\end{tabular}
\caption{\textbf{Comparison to Other Platforms}, we illustrate the difference between our benchmark and other popular platforms for studying multi-agent coordination or embodied agents. 
\textbf{Multi-turn communication} refers to the ability of agents to ask follow up questions and engage in a grounded dialogue. 
\textbf{Partial Observability} refers to agents not being being fully aware of everything the other agent perceives, as this is a necessity for testing Theory of Mind capabilities. \textbf{Long horizon} refers to the complex sequence of actions (on average over 20 steps) that need to be taken in order to accomplish our task objectives. \textbf{Quantitative Evaluation} refers to the capability of the tasks to be evaluated for collaboration, which is done qualitatively in previous papers\cite{park2023generative}.
}
\label{tab:simulator_differences}
\end{table*}
